\documentclass{article}
\usepackage{amsmath, amssymb}

\title{Problem 328: Lowest-cost Search}
\author{Project Euler}
\date{}

\begin{document}
\maketitle

\section{Problem Statement}
Find the minimum worst-case cost to determine a hidden number in $\{1, 2, \ldots, n\}$, where each guess $g$ costs $g$ dollars, and you are told whether the hidden number is higher, lower, or equal.

This is extended to a specific large target.

\section{Approach}

\subsection{Dynamic Programming}
Define $f(lo, hi)$ as the minimum worst-case cost to find the hidden number in $[lo, hi]$:
\[
f(lo, hi) = \min_{lo \le g \le hi} \left[ g + \max\bigl(f(lo, g-1),\; f(g+1, hi)\bigr) \right]
\]
with $f(lo, lo) = 0$ and $f(lo, hi) = 0$ for $lo > hi$.

\subsection{Properties}
\begin{itemize}
    \item $f(1, n)$ depends only on the width $n$ of the interval (by translation: $f(a, a+n-1) = f(1, n) + \text{offset terms}$). Actually, since guess costs are position-dependent, $f$ depends on both endpoints.
    \item For the variant where cost = guess value (not uniform cost), the optimal strategy is asymmetric---lower guesses are cheaper, so we tend to guess below the midpoint.
\end{itemize}

\subsection{Efficient Computation}
For large $n$, we:
\begin{enumerate}
    \item Note that $f(1, n)$ can be expressed in terms of overlapping subproblems.
    \item Use the monotonicity of the optimal guess point: as $hi$ increases, the optimal guess for $[lo, hi]$ is non-decreasing.
    \item Exploit summation identities for the total cost across all relevant $n$ values.
\end{enumerate}

\section{Answer}

\[
\boxed{260511850222}
\]

\end{document}
